\documentclass[11pt,a4paper]{article}
\usepackage[T1]{fontenc}
\usepackage{lmodern}
\usepackage[margin=20mm]{geometry}
\usepackage{amsmath,amssymb,amsthm}
\usepackage{microtype}
\usepackage[hidelinks]{hyperref}
\hypersetup{
  pdftitle={Erdos 252: proof outline},
  pdfauthor={Tokengrinder},
  pdfsubject={Irrationality of the factorial divisor-sum series},
  pdfkeywords={Erdos 252, irrationality, Lean, proof outline}
}
\newtheorem*{theorem}{Theorem}
\newcommand{\N}{\mathbb{N}}
\newcommand{\Z}{\mathbb{Z}}
\newcommand{\R}{\mathbb{R}}
\newcommand{\stirling}[2]{\left\{\begin{smallmatrix}#1\\#2\end{smallmatrix}\right\}}
\setlength{\parindent}{0pt}
\setlength{\parskip}{2pt plus 1pt}
\linespread{0.94}
\AtBeginDocument{%
  \setlength{\abovedisplayskip}{6pt plus 1pt minus 1pt}%
  \setlength{\belowdisplayskip}{6pt plus 1pt minus 1pt}%
  \setlength{\abovedisplayshortskip}{0pt plus 1pt}%
  \setlength{\belowdisplayshortskip}{4pt plus 1pt minus 1pt}%
}
\setlength{\emergencystretch}{2em}
\makeatletter
\renewcommand\section{\@startsection{section}{1}{\z@}%
  {-1.7ex plus -.3ex minus -.2ex}{.8ex plus .2ex}%
  {\normalfont\large\bfseries}}
\renewcommand\paragraph{\@startsection{paragraph}{4}{\z@}%
  {1.2ex plus .3ex minus .2ex}{-1em}%
  {\normalfont\normalsize\bfseries}}
\makeatother
\pagestyle{plain}

\begin{document}

\begin{center}
  {\Large\bfseries Erd\H{o}s 252}\par
  \smallskip
  {\large Irrationality of the factorial divisor-sum series}\par
  \medskip
  Tokengrinder
\end{center}

\textbf{Proof outline.}
This is a mathematical exposition of the accompanying Lean
formalization. It describes the main constructions and the contradiction;
the complete rigorous proof, including all estimates and finite arithmetic,
is the Lean development linked below.

For a positive integer $n$ and $k\in\N=\{0,1,2,\ldots\}$, put
\[
  \sigma_k(n)=\sum_{d\mid n}d^k,
  \qquad
  \alpha_k=\sum_{n=1}^{\infty}\frac{\sigma_k(n)}{n!}.
\]
The series converges for every $k$. In the formal statement the sum starts
at $n=0$; the corresponding divisor-sum term is zero.

\begin{theorem}
For every $k\in\N$, the real number $\alpha_k$ is irrational.
\end{theorem}

\section*{Positive exponents}

Fix $k\geq 1$ and suppose, for a contradiction, that $\alpha_k$ is rational.
Clearing its denominator by a sufficiently large factorial shows that
\[
  T_n=n!\sum_{m>n}\frac{\sigma_k(m)}{m!}\in\Z
  \qquad\text{for all sufficiently large }n.
\]
All grids, coefficients, shifts, and moduli below are fixed once $k$ is fixed.

\paragraph{1. Expand the tails.}
Let $\stirling{a}{b}$ denote a Stirling number of the second kind.
Expanding the first $k+1$ factorial denominators in reciprocal powers gives
\[
  T_n=
  \sum_{h=1}^{k+1}\sum_{\ell=1}^{k+1}
    \stirling{\ell-1}{h-1}
    \frac{\sigma_k(n+h)}{(n+h)^\ell}
  +O_k(n^{-3/2}).
\]
Pairing each divisor with its complementary divisor bounds the divisor
count by a constant times $\sqrt n$. Consequently
$\sigma_k(n)=O_k(n^{k+1/2})$. Together with a geometric bound for the
omitted tail, this controls both sources of expansion error.

\paragraph{2. Align shifted arguments.}
Use the finite grid $E=\{0,1,\ldots,k+1\}^k$. One explicit choice in the
formalization is
\[
  b=k+2,\qquad F=b^k-1,\qquad D=F!,\qquad B=1+kDF,
\]
\[
  p_e=B+D\sum_{j=0}^{k-1}e_jb^j,
  \qquad
  t_e=D\sum_{j=0}^{k-1}(j+1)e_jb^j
  \qquad(e\in E).
\]
The integers $p_e$ are pairwise coprime, and the shifts
$r_{e,h}=hp_e-t_e$ are positive for $1\leq h\leq k+1$.
The Chinese remainder theorem supplies a progression $N=A+Qm$, $A,m\in\N$, with
$Q=\prod_{e\in E}p_e^2$, satisfying $N\equiv t_e\pmod{p_e^2}$ for every $e$.
For sufficiently large $N$, the integer $n_e=(N-t_e)/p_e$ is positive.
The squared congruence gives $p_e\mid n_e$. Every retained $h$ divides $D$,
while $\gcd(p_e,D)=1$, so $\gcd(p_e,n_e+h)=1$. Multiplicativity gives
\[
  \sigma_k(p_e)\sigma_k(n_e+h)
    =\sigma_k(N+r_{e,h})
  \qquad(1\leq h\leq k+1).
\]
None of the multipliers $p_e$ is required to be prime.

\paragraph{3. Cancel the large terms.}
Assign the integer binomial weights
\[
  w_e=\prod_{j=0}^{k-1}(-1)^{k+1-e_j}\binom{k+1}{e_j}.
\]
For a term with $1\leq h\leq\ell\leq k$, varying coordinate $e_{h-1}$
leaves the shift $r_{e,h}$ fixed. The remaining multiplier is a polynomial
of degree at most $k$ in that coordinate. Its order-$(k+1)$ finite
difference vanishes. Terms with $h>\ell$ have zero Stirling coefficient.
Thus every denominator order $\ell\leq k$ cancels exactly.

To name the surviving terms, let
\[
  I=E\times\{1,\ldots,k+1\},\qquad
  c_{e,h}=w_ep_e^{k+1}\stirling{k}{h-1},\qquad
  \rho_k(x)=\frac{\sigma_k(x)}{x^k}\quad(x\in\N,\ x\geq1).
\]
The weighted tail $\sum_e w_e\sigma_k(p_e)T_{n_e}$ is an integer for all
sufficiently large $N$ in the progression. After cancellation it equals
\[
  \sum_{(e,h)\in I}
    \frac{c_{e,h}\rho_k(N+r_{e,h})}{N+r_{e,h}}
    +O_k(N^{-3/2}).
\]
Since $\rho_k(x)=O_k(\sqrt x)$, this integer tends to zero and hence is
eventually zero. Multiplying by $N$, the error still tends to zero.
Replacing $N/(N+r_{e,h})$ by $1$ introduces only another vanishing error.
It follows that
\[
  R(N):=\sum_{(e,h)\in I}c_{e,h}\rho_k(N+r_{e,h})
  \longrightarrow0
  \quad\text{along }N=A+Qm.
\]

\paragraph{4. Contradict the remaining limit.}
The shift $r_*=(k+1)B$, indexed by the zero vertex and $h=k+1$, occurs
exactly once; its coefficient $c_*$ is nonzero. For a positive step $Q$
and positive offset $a$, the actual arithmetic-progression mean is
\[
  \lim_{M\to\infty}\frac1M\sum_{j=0}^{M-1}\rho_k(Qj+a)
   =\mu_Q(a)
   :=\sum_{\substack{d\geq1\\\gcd(d,Q)\mid a}}
       \frac{\gcd(d,Q)}{d^{k+1}}>0.
\]
This includes $k=1$: summable domination is applied to averaged divisor
counts, not to an unjustified uniform bound on the pointwise divisor sum.

Choose a fresh prime $L$ larger than $Q$ and every shift. Refining the
progression by a residue $v\in\{0,\ldots,L-1\}$ changes its individual means by the factor
\[
  \mu_{QL}(Qv+a)
   =\mu_Q(a)\left(1-L^{-(k+1)}
          +\mathbf{1}_{\{L\mid Qv+a\}}L^{-k}\right),
\]
where $\mathbf{1}$ is the indicator of the stated divisibility condition.
One residue makes $Qv+A$ divisible by $L$ and therefore misses every
positive shift. Another makes $Qv+A+r_*$ divisible by $L$ and hits only
the unique shift $r_*$. The two resulting means of $R$ differ by
\[
  c_*\,\mu_Q(A+r_*)\,L^{-k}\ne0.
\]
But both means must be zero if $R(A+Qm)\to0$, because the same limit holds
on either subprogression and under Ces\`aro averaging. This is the desired
contradiction.

\section*{Exponent zero}

For $k=0$, divisor pairing and a geometric estimate directly give
$0<T_n=O(n^{-1/2})$. These positive tails cannot be eventually integral:
for large $n$ they lie strictly between $0$ and $1$. Rationality would
force precisely that eventual integrality, so $\alpha_0$ is irrational too.

\section*{Formal verification and sources}

\begingroup\small
The complete theorem is \texttt{Erdos252.erdos\_252}. Its checked axiom
dependencies are exactly \texttt{propext}, \texttt{Classical.choice}, and
\texttt{Quot.sound}. Build, statement audits, and fresh replay using
Lean's kernel have passed with Lean 4.33.1 and the pinned mathlib revision.
Fresh replay uses Lean's own kernel, not a separate kernel implementation.
No prime-pattern conjecture, distribution assumption, or finite
computational surrogate is used.

The \href{https://github.com/tokengr1nder/Erdos252}{source repository}
provides reproduction instructions and exact dependency revisions. Its
formal entrypoints are the
\href{https://github.com/tokengr1nder/Erdos252/blob/main/Erdos252/Solution.lean}{complete theorem},
the \href{https://github.com/tokengr1nder/Erdos252/blob/main/Erdos252/DilationIrrationality.lean}{positive-exponent proof},
the \href{https://github.com/tokengr1nder/Erdos252/blob/main/Erdos252/DilationZeroCase.lean}{zero-exponent proof},
and the \href{https://github.com/tokengr1nder/Erdos252/blob/main/audit/Statement.lean}{statement audit}
(definitions, summability, and reindexing). See also the
\href{https://www.erdosproblems.com/252}{original problem statement}.
\par\endgroup

\end{document}
